% =====================================================================
% EdgeFlow 用户手册 · 第 9 章 常见问题与故障排查
% 事实依据：docs/KEADM.md §5、docs/REAL-CLUSTER-GUIDE.md、
%           docs/EVENTBUS-GUIDE.md、docs/API-SPEC.md
% =====================================================================
\chapter{常见问题与故障排查}
\label{ch:9}

本章汇集 v0.7.0 实测与文档沉淀的常见问题。排障第一步永远是\textbf{看日志}，
先定位故障在哪一层：云边通道、容器运行时、设备数据面，还是 API 层。

\section{日志位置}

\begin{table}[htbp]
\centering
\caption{各组件日志位置}
\label{tab:9-logs}
\begin{tabularx}{\textwidth}{l X}
\toprule
\textbf{组件} & \textbf{日志位置} \\
\midrule
cloudcore & stdout / 日志文件（部署形态而定） \\
edgecore（systemd） & \texttt{journalctl -u edgecore -e} / \texttt{-f} 跟踪 \\
edgecore（裸进程） & stdout \\
审计台账 & \texttt{data/audit-ledger.jsonl}（云端） \\
操作台账 & \texttt{<output-dir>/ops-ledger.jsonl}（keadm） \\
\bottomrule
\end{tabularx}
\end{table}

\section{常见问题速查表}

\begin{table}[htbp]
\centering
\caption{常见问题与处置}
\label{tab:9-faq}
\begin{tabularx}{\textwidth}{l l X}
\toprule
\textbf{现象} & \textbf{常见原因} & \textbf{处置} \\
\midrule
edgecore 起不来 & 网络不通 / 配置错误 &
\texttt{journalctl -u edgecore -e} 查看；确认可达
\texttt{--cloudcore-ip} 的 hub 端口；确认 env 文件键值未被手改 \\
EventBus 连接失败 & MQTT broker 未启动 / 地址错 &
查看日志确认已\textbf{降级本地模式}（边缘照常运行）；确认 broker
监听 1883；macOS 注意 broker 二进制在
\texttt{/opt/homebrew/sbin/mosquitto} \\
MQTT 建连超时 & broker 无响应 / 网络隔离 &
调整 \texttt{EDGEFLOW\_EDGECORE\_MQTT\_CONNECT\_TIMEOUT}（默认 5s）；
确认 \texttt{EDGEFLOW\_EDGECORE\_MQTT\_ADDR} \\
容器不启动 & Docker 未安装 / daemon 未启动 &
Edged 依赖 Docker：确认 \texttt{docker info} 正常；daemon 启动后
Edged 下轮调谐（默认 5s）自动拉起 \\
端口冲突 & 8080 / 10000 / 1883 被占用 &
\texttt{lsof -i :8080} 等定位占用进程；改用
\texttt{EDGEFLOW\_CLOUDCORE\_PORT} / \texttt{HUB\_PORT} /
\texttt{EDGEFLOW\_EDGECORE\_MQTT\_ADDR} \\
节点未注册 404 & 节点从未注册 / 已离线 &
检查 edgecore 是否在跑、注册是否成功；\texttt{GET /api/v1/nodes}
确认节点状态 \\
下发 504 超时 & 边缘宕机 / 链路抖动 &
确认边缘进程存活；504 可重试，边缘侧幂等去重；持续 504 检查
网络与心跳 \\
容器反复重启 & 业务崩溃 / 退出型镜像 &
观察 \texttt{RestartCount} 与 CrashLoopBackOff（3 次阈值/30s 退避/
60s 稳定重置）；查看容器日志定位业务根因 \\
云端重启后节点消失 & v0.4.0 起云端\textbf{分级持久化} &
节点\textbf{台账保留}，重启后短暂 \texttt{Unknown}，边缘 \textbf{≤1 个
心跳周期}自动翻新 \texttt{Ready}（无需重新注册）；Pod/设备上报数据
重启后 ≤1 个上报周期自愈；长时间未恢复再查日志 \\
全量软离线（外部多副本） & etcd 集群故障超过 \texttt{NODE\_LEASE\_TTL} &
外部 etcd 多副本形态的\textbf{预期行为}：节点短时 \texttt{Offline}，
etcd 恢复后 \textbf{≤1 个心跳周期自愈、零数据删除}；监控告警阈值按
≈2×TTL 折算，避免误报 \\
日志出现 \texttt{concurrent-write} & 多副本并发写同一设备 Desired &
CAS 冲突重试（≤3 次）耗尽属\textbf{正常语义}：HTTP 仍 \texttt{200}，
日志提示并发写冲突；重新下发指令收敛即可 \\
模型发布返回 409/422 & 同模型在途发布 / 目标版本非 active /
无 Ready 节点 / 白名单含未知节点 / 回滚被新版本接管 &
409 = 冲突（读 error 文案，在途发布含 releaseID）；422 = 业务前置
不满足（响应含 \texttt{unknownNodes} 等上下文）；修正后重试 \\
外部模式启动失败 & 探活失败 / 鉴权被拒 / 明文护栏 &
日志区分 \texttt{Unavailable}（集群不可达）与 \texttt{PermissionDenied}
（凭据/角色不足）；非回环 + 无 TLS 被护栏拒绝时配置 TLS 或内网逃生门；
处理见部署指南 §10.7.6 \\
纯内存模式发布任务重启丢失 & 发布/部署影子为内存态（L22） &
纯内存模式（\texttt{ETCD\_ENABLED=false}）下模型发布任务与部署影子
重启丢失；生产建议使用 embed/外部 etcd 模式 \\
\bottomrule
\end{tabularx}
\end{table}

\section{分层排查指引}

\subsection{云边通道层}

\begin{enumerate}
  \item 确认 edgecore 进程存活且已注册：
\begin{lstlisting}
curl http://127.0.0.1:8080/api/v1/nodes
\end{lstlisting}
  \item 节点存在但状态 \texttt{Offline}：心跳停滞超过 180s（约 6 个
        30s 心跳周期；\textbf{外部多副本形态}判活为 etcd 租约，阈值取
        \texttt{EDGEFLOW\_CLOUDCORE\_NODE\_LEASE\_TTL}（默认 300s），
        见第 5 章）。检查网络、云端进程；
  \item 节点不存在（404）：从未注册或注册失败。看 edgecore 日志中
        Register/RegisterAck 是否成功；
  \item 启用 mTLS 后连不上：先确认两侧开关都开、证书目录正确、
        边缘连接地址在服务端证书 SAN 内（见第~\ref{ch:8}~章）。
\end{enumerate}

\subsection{容器运行层}

\begin{enumerate}
  \item 调谐周期默认 5s，改动 Pod 定义后最多等一个周期生效；
  \item 容器反复拉起并进入退避：区分「业务崩溃」与「退出型镜像」。
        CrashLoopBackOff 是保护机制，不是故障本身；
  \item 镜像漂移检测会在期望镜像与运行镜像不一致时自动重建（每轮 1 个，
        逐轮滚动），滚动期间短暂重启属预期。
\end{enumerate}

\subsection{设备数据面层}

\begin{enumerate}
  \item \texttt{GET /api/v1/devices} 看不到设备：确认 edgecore 内
        Mapper 已注册（日志有「Mapper 注册完成」），MQTT 模式下确认
        broker 可达；
  \item 设备属性不更新：确认采集周期（mock\_sensor 默认 2s）与
        上报周期（默认 30s）；
  \item 指令下发失败：按状态码区分——404 未送达（节点离线，修通道）、
        502 送达被拒（改请求）、504 未确认（重试）。
\end{enumerate}

\section{API 错误语义速记}

\begin{itemize}
  \item \texttt{404}：节点未注册/离线，\textbf{无需重试}，先修通道；
  \item \texttt{502}：消息已送达但边缘拒绝，\textbf{重试无意义}，
        检查请求内容；
  \item \texttt{504}：可能送达但未确认，\textbf{可重试}，边缘幂等去重；
  \item 错误响应统一为 JSON：\texttt{\{"error":"<机器可读原因>"\}}。
\end{itemize}

完整语义见附录~B。

\section{环境相关的已知边界}

\begin{itemize}
  \item \textbf{macOS/Docker Desktop}：宿主机无法直接路由 kind 节点 IP
        的 NodePort（实测 connection reset），用
        \texttt{kubectl port-forward} 替代；Linux 上可直接
        \texttt{nodeIP:nodePort} 访问；
  \item \textbf{真实多节点集群}：v0.1.0 验证范围为 kind 单节点集群 +
        单边缘节点（2026-08-14 实测），多节点（10+）E2E 与压测未做；
  \item \textbf{镜像仓库}：镜像未推送远端仓库，跨机部署需自建 registry
        （\texttt{docker run -d -p 127.0.0.1:5001:5000 registry:2} +
        buildx 双架构构建，见 docs/MULTIARCH.md）；
  \item \textbf{nodeID 字符约束}（v0.4.0 起）：节点 ID 必须匹配正则
        \texttt{\textasciicircum{}[A-Za-z0-9.\_\-]+\$}，含 \texttt{/} 的
        nodeID 拒绝写入云端台账。
\end{itemize}
